% Online sync refresh: revised workshop report, 2026-08-29.
The definition of sublinear regret requires
\[
    \lim_{T\to\infty}
    \frac{\operatorname{Regret}_T}{T}
    =
    0,
\]
which is the standard asymptotic criterion used in multi-armed bandit
regret analysis. Therefore, a natural numerical
test is to estimate the regret at a sequence of increasingly large horizons
and examine its asymptotic growth.

Suppose the experiment is repeated independently for $M$ runs, and let
$A_t^{(r)}$ denote the arm selected at time $t$ in run $r$. Since the
synthetic bandit environment is known to us, the optimal expected reward is
\[
    R^\ast
    =
    \max_{a\in\mathcal A} q^\ast(a).
\]
The true action values $q^\ast(a)$ are used only for performance evaluation
after an action has been selected; they are not made available to the
learning algorithm.

For run $r$, define the cumulative pseudo-regret at horizon $T$ as
\[
    C_T^{(r)}
    =
    \sum_{t=1}^{T}
    \left[
        R^\ast-q^\ast\!\left(A_t^{(r)}\right)
    \right].
\]
Averaging over the $M$ independent reruns gives the empirical estimate
\[
    \widehat{\operatorname{Regret}}_T
    =
    \frac{1}{M}
    \sum_{r=1}^{M}
    C_T^{(r)}.
\]

The first diagnostic is the normalised empirical regret
\begin{equation}
    \widehat{\rho}_T
    =
    \frac{\widehat{\operatorname{Regret}}_T}{T}.
    \label{eq:q2-normalised-regret}
\end{equation}
By the definition of sublinear regret,
\eqref{eq:q2-normalised-regret} should approach zero as $T$ increases.
However, examining this quantity at only one finite horizon can be
misleading. For example, if
\[
    \operatorname{Regret}_T=C+cT,
    \qquad c>0,
\]
then
\[
    \frac{\operatorname{Regret}_T}{T}
    =
    \frac{C}{T}+c
\]
may decrease substantially with $T$, even though it eventually converges
to the positive constant $c$ and the regret is therefore linear.

To distinguish such transient behaviour from genuinely sublinear growth,
we additionally evaluate the incremental regret rate between successive
horizons $T_i<T_{i+1}$:
\begin{equation}
    \widehat{s}_i
    =
    \frac{
        \widehat{\operatorname{Regret}}_{T_{i+1}}
        -
        \widehat{\operatorname{Regret}}_{T_i}
    }{
        T_{i+1}-T_i
    }.
    \label{eq:q2-incremental-rate}
\end{equation}
This is the average slope of the cumulative-regret curve over the interval
$[T_i,T_{i+1}]$.

The use of \eqref{eq:q2-incremental-rate} follows directly from the
definition of sublinear regret. Consider geometrically increasing horizons
$T_{i+1}=cT_i$, where $c>1$. Since cumulative pseudo-regret is non-negative
and non-decreasing,
\[
    0
    \leq
    \widehat{s}_i
    \leq
    \frac{
        \widehat{\operatorname{Regret}}_{cT_i}
    }{
        (c-1)T_i
    }
    =
    \frac{c}{c-1}
    \frac{
        \widehat{\operatorname{Regret}}_{cT_i}
    }{
        cT_i
    }.
\]
Hence, if the underlying regret is sublinear, the right-hand side approaches
zero as the horizon increases, and therefore the incremental regret rate
should also approach zero. In contrast, if
\[
    \operatorname{Regret}_T
    \sim c_0 T,
    \qquad c_0>0,
\]
then
\[
    \widehat{s}_i
    \longrightarrow
    c_0,
\]
so a positive late-horizon plateau in the incremental regret rate provides
numerical evidence of linear regret.

As a supplementary diagnostic, the local effective growth exponent is
estimated by
\[
    \widehat{\alpha}_i
    =
    \frac{
        \log
        \widehat{\operatorname{Regret}}_{T_{i+1}}
        -
        \log
        \widehat{\operatorname{Regret}}_{T_i}
    }{
        \log T_{i+1}-\log T_i
    }.
\]
If locally
$\widehat{\operatorname{Regret}}_T\approx CT^\alpha$, then
$\widehat{\alpha}_i\approx\alpha$. Thus,
$\widehat{\alpha}_i$ approaching $1$ supports locally approximately linear
power-law growth, whereas values remaining below $1$ support a sublinear
power-law interpretation. However, $\widehat{\alpha}_i$ is neither a
necessary nor a sufficient test for sublinear regret. For example, the
sublinear function $T/\log T$ also has a local effective exponent approaching
$1$. The normalised regret and incremental regret rate are therefore the
primary diagnostics, while the effective exponent is supplementary.

The proposed numerical procedure is therefore to simulate each algorithm
up to a maximum horizon and record its cumulative pseudo-regret at several
geometrically increasing horizons, for example
\[
    T\in
    \{10^3,\,3\times10^3,\,10^4,\,3\times10^4,\,10^5\}.
\]
The experiment is repeated for many independent runs and
$\widehat{\operatorname{Regret}}_T$,
$\widehat{\rho}_T$,
$\widehat{s}_i$, and
$\widehat{\alpha}_i$
are calculated.

Numerical behaviour is considered consistent with sublinear regret when,
over the largest simulated horizons,
\[
    \widehat{\rho}_T
    \longrightarrow 0
    \qquad\text{and}\qquad
    \widehat{s}_i
    \longrightarrow 0,
\]
based on their observed multi-horizon trends. Conversely, a positive
late-horizon incremental regret rate provides evidence of a linear
component. The effective exponent may support a local power-law
interpretation, but it is not used as an automatic classification rule.

Because only finitely many horizons and reruns can be simulated, this test
cannot prove an asymptotic regret property. It provides numerical evidence
by examining whether the observed large-horizon growth is consistent with
the mathematical definition of sublinear regret.


The proposed numerical procedure is summarised in
Algorithm~\ref{alg:q2-sublinear-test}.

\begin{algorithm}[H]
\caption{Multi-horizon numerical test for sublinear regret}
\label{alg:q2-sublinear-test}
\begin{algorithmic}[1]
\Require Bandit policy $\pi$, horizons
$T_1<\cdots<T_L$, number of reruns $M$

\For{$r=1,\ldots,M$}
    \State Simulate $\pi$ up to $T_L$
    \State Record cumulative pseudo-regret
    $C_{T_i}^{(r)}$ at each horizon $T_i$
\EndFor

\For{$i=1,\ldots,L$}
    \State Estimate
    $\widehat{\operatorname{Regret}}_{T_i}
    \gets \frac{1}{M}\sum_{r=1}^{M}C_{T_i}^{(r)}$
    \State Compute
    $\widehat{\rho}_{T_i}
    \gets \widehat{\operatorname{Regret}}_{T_i}/T_i$
\EndFor

\For{$i=1,\ldots,L-1$}
    \State Compute incremental regret rate
    \[
        \widehat{s}_i
        \gets
        \frac{
            \widehat{\operatorname{Regret}}_{T_{i+1}}
            -
            \widehat{\operatorname{Regret}}_{T_i}
        }{
            T_{i+1}-T_i
        }
    \]
    \State Compute effective exponent
    \[
        \widehat{\alpha}_i
        \gets
        \frac{
            \log \widehat{\operatorname{Regret}}_{T_{i+1}}
            -
            \log \widehat{\operatorname{Regret}}_{T_i}
        }{
            \log T_{i+1}-\log T_i
        }
    \]
\EndFor

\State Examine the late-horizon trends of the primary diagnostics
$\widehat{\rho}_{T_i}$ and $\widehat{s}_i$

\State Use $\widehat{\alpha}_i$ only as a supplementary local
power-law diagnostic

\State Interpret the measured trends manually; do not apply a fixed
numerical classification threshold

\end{algorithmic}
\end{algorithm}
